% problem-000080 4 化学反应与结构综合分析
\section*{4. 化学反应与结构综合分析}
制备⽯硫合剂（⼀种常⽤农药）的⽅法如下：先在反应器内加⽔使⽯灰消解，然后加⾜⽔量，在搅拌下把硫磺粉慢慢倒⼊。升温熬煮，使硫发⽣歧化反应。滤去固体渣滓得到粘稠状深棕⾊透明碱性液体成品。

\subsection*{4-1}
写出⽣成产品的歧化反应的化学反应⽅程 $\updownarrow \updownarrow _ { \updownarrow }$

\subsection*{4-2}
说明固体渣滓的成分（假定⽔和⽣⽯灰皆为纯品）。

\subsection*{4-3}
写出呈现深棕⾊物质的化学式。

\subsection*{4-4}
通常采⽤下述路线进⾏⽯硫合剂成品中有效活性物种的含量分析。写出A、B和C所代表的物种的化学式。

（图：）

\subsection*{4-5}
我国科技⼯作者⾸创⽤⽯硫合剂提取⻩⾦的新⼯艺，其原理类似于氰化法提⾦，不同的是过程中不需要通⼊空⽓。写出⽯硫合剂浸⾦过程中氧化剂的化学式，指出形成的Au(I)配离⼦中的配位原⼦。
